\documentclass[10pt,a4paper]{amsart}
\usepackage[margin=19mm]{geometry}
\usepackage{amsmath,amssymb,mathtools}
\usepackage[scr=rsfs,cal=euler]{mathalfa}
\usepackage{xfrac}
\usepackage[unicode,hidelinks,bookmarks=false]{hyperref}
\newcommand{\ie}{i.e.\ }
\newcommand{\cf}{cf.\ }
\newcommand{\resp}{respectively\ }
\newcommand{\Subsection}{\subsection}
\providecommand{\nobreakdash}{\nobreak\hbox{-}}
\setlength{\emergencystretch}{2em}
\multlinegap0pt
\usepackage[shortlabels]{enumitem}
\usepackage{amssymb}
\usepackage{mathtools}
\newtheorem{lemma}{Lemma}
\newcommand{\Stab}[2]{\operatorname{Stab}_{#1}\left(#2\right)}
\renewcommand{\phi}{\varphi} 
\title{A combination theorem for the twist conjecture for Artin groups}
\author{Oli Jones, Giorgio Mangioni, Giovanni Sartori}
\date{Source: arXiv:2507.13971v2 (CC BY 4.0)}
\begin{document}
\maketitle

\noindent\emph{Source and licence.} A combination theorem for the twist conjecture for Artin groups. Version arXiv:2507.13971v2 (CC BY 4.0), distributed by arXiv under the Creative Commons Attribution 4.0 International licence, http://creativecommons.org/licenses/by/4.0/, as recorded in the arXiv licence metadata for this exact version and shown on the abstract page https://arxiv.org/abs/2507.13971v2. Full licence notice: \texttt{licenses/arxiv-2507.13971-CC-BY-4.0.txt}.\par\smallskip

\noindent\textit{Editorial scope.} Counted unit: the article's lemma lem:stabsOfCollapse (source0.tex lines 299--307). The article's complete original proof of it is delivered with it (source0.tex lines 309--321, one \texttt{proof} environment of 13 lines). No mathematics is written, repaired or re-derived: the statement and the proof are the article's own lines, in the article's own notation, and no retained line is substituted or re-flowed.  No further source-local context is needed: every cross-reference of every retained line resolves inside the two delivered spans.  Nothing was omitted from any retained span, so the package carries no omission record.  No retained line contains a citation, so the wrapper carries no bibliography and no bibliography entry.  No retained line contains a figure environment, an image-inclusion command or drawing code, so no image is delivered and the package declares no asset dependency.  Editorial formatting only, in addition: an AMS \texttt{amsart} preamble that restates the article's own declarations and macros used by the retained lines, namely the environments \texttt{lemma} on separate counters, together with the standard packages those lines need.  The retained environments print in this document as Lemma 1 in order of appearance, whereas the article prints the same items with the numbering of its own full document.  The complete native source of the article as distributed in the arXiv e-print for this exact version is bundled unchanged as \texttt{sources/181817/source0.tex}.
\begin{lemma}
    \label{lem:stabsOfCollapse}
    Let $(T,\Omega)$ be a $G$-tree, let $e=\{u,v\}$ be an edge of $T$ with $u,v$ in different $\Omega$-orbits and $\Stab \Omega u\le\Stab\Omega v$, and let $\phi\colon(T,\Omega)\to(T',\Omega')$ be the elementary collapse of the edge~$e$. Every edge stabiliser (resp. vertex stabiliser) for $\Omega'$ is also an edge stabiliser (resp. vertex stabiliser) for $\Omega$.
     In particular:
    \begin{enumerate}
        \item if $x'$ is an edge of $T'$ (resp. a vertex of $T'$ not in the orbit of $\phi(e)$), then for every edge (resp. vertex) $x$ of $T$ such that $\phi(x) = x'$, $\Stab\Omega x=\Stab{\Omega'}{x'}$;
        \item $\Stab{\Omega'}{\phi(e)}=\Stab\Omega v$.
    \end{enumerate}
\end{lemma}

\begin{proof}
We first notice that, if $x'$ is an open edge of $T'$ (note that restricting to the interior of edges does not change the stabilisers), or a vertex which is not in the $G$-orbit of $\phi(e)$, then $\phi$ is injective on $\phi^{-1}(x')$. Now suppose $\phi(x) = x'$. Then if $gx' = x'$ it follows by $G$-equivariance that $\phi(x)= x'= gx' = \phi(gx)$, so $x = gx$ by injectivity. Conversely if $gx = x$ then $gx' = g\phi(x) = \phi(gx) = x'$ by equivariance. Hence, we have that $\Stab{\Omega'}{x'}=\Stab{\Omega}{x}$ for any edge (resp. vertex) $x$ such that $\phi(x) = x'$.

    
To conclude the proof, we now show that
\begin{equation}\label{eq:stab_phi(e)}
    \Stab{\Omega'}{\phi(e)}=\Stab{\Omega}{v}.
\end{equation}
    For the inclusion $\le$, let $g\in \Stab{\Omega'}{\phi(e)}$. Towards a contradiction, assume that $g\cdot v\neq v$. Because $g$ fixes $\phi(e)$, the vertices $v$ and $g\cdot v$ have the same image under~$\phi$. Then there is a path of $G$-translates $e_1=g_1\cdot e, e_2=g_2\cdot e_1,\ldots, e_k=g_k \cdot e_{k-1}$ of $e$ connecting $v$ to $g\cdot v$. Notice that $k\ge 2$, as $v$ and $u$ are in different $G$-orbits; moreover $g_1\cdot v=v$, as the path connects $v$ to $g\cdot v$, and for the same reason $g_2$ fixes $g_1 \cdot u$ but not $v$. However $g_2 \in \Stab{\Omega}{g_1 \cdot u}=\Stab{\Omega}{u}^{g_1}\le \Stab{\Omega}{v}$, which is absurd.

    For the converse inclusion, if $g\in \Stab{\Omega}{v}$ then $$g\phi(e)=g\phi(v)=\phi(g\cdot v)=\phi(v)=\phi(e),$$
    where we used that $\phi$ is $G$-equivariant.
\end{proof}

\end{document}
